\documentclass[11pt]{article}
\usepackage[margin=1in]{geometry}
\usepackage{times}
\usepackage{hyperref}
\hypersetup{colorlinks=true, linkcolor=black, urlcolor=blue, citecolor=black}
\title{Preorders Are Not a Product: The Perendev Magnetic Motor and the Financing-First Pattern}
\author{Flyxion \\ \small Independent Researcher}
\date{}
\begin{document}
\maketitle

\begin{abstract}
The Perendev magnetic motor, promoted by Mike Brady's company of the same name from the mid-2000s, claimed to use a specific arrangement of permanent magnets to produce continuous rotation and net usable power with no external energy input, and was marketed internationally through licensing agreements and direct unit preorders before any independently verified working unit had been demonstrated. This paper argues that the device fails on the same permanent-magnet energy conservation grounds addressed elsewhere in this series, and treats the case additionally as an example of a financing structure, licensing and preorder revenue collected ahead of independent product verification, that generates a strong incentive to maintain a claim regardless of whether the underlying device works, independent of the inventor's original intentions.
\end{abstract}

\section{Introduction}

Mike Brady's Perendev company claimed, from around 2005 onward, to have developed a magnetic motor generating rotational output from a fixed arrangement of static permanent magnets on a rotor and housing, requiring no external power source once started, and pursued a business model of selling regional licensing rights and taking direct preorders for motor units marketed for applications including home power generation and electric vehicle propulsion. No independently verified working unit was ever delivered to a customer or demonstrated under conditions ruling out an external power source, and Brady was later convicted, in a German court in 2011, of fraud related to the collection of funds for undelivered motors.

\section{The Physics Objection Restated}

As with the Yildiz Motor addressed elsewhere in this series, a rotor driven solely by the interaction of static permanent magnets does no net work over a closed rotational cycle, because the magnetic field of an unpowered permanent magnet is fixed and conservative in the relevant sense; any energy a passing magnet appears to supply during one part of a rotor's cycle must be returned, less dissipative losses, during another part of the same cycle, a consequence of energy conservation applied to a magnetostatic field with no external power replenishing it. The claimed design, magnets arranged asymmetrically or shielded at specific points in the rotation to allegedly prevent the usual retarding force from balancing the driving force, describes an engineering approach to timing a conservative force's application, not a mechanism that evades the conservation law itself, and no published account of the Perendev design specifies a shielding or timing mechanism that has been shown, independently of the inventor's own claims, to produce net output exceeding input under controlled measurement.

\section{Why the Financing Structure Deserves Separate Attention}

Independent of the underlying physics, the Perendev business model, collecting licensing fees and unit preorders substantially in advance of any independently verified working product, creates a direct financial incentive structure in which continued claims of imminent commercial delivery generate revenue regardless of whether a working unit exists, and in which admitting the device does not work carries an immediate and severe financial and legal cost to whoever has already collected on that basis. This does not, by itself, establish that Brady set out to defraud rather than genuinely believing in an unworkable design and continuing to promote it past the point of correction, but it does mean that the persistence of promotional claims well past the point of any independent verification is not, in this case, surprising or informative about the underlying physics one way or the other, since the financial incentive to continue promoting exists regardless of whether the device works.

\section{The Legal Resolution}

The 2011 German fraud conviction is one of relatively few cases in this genre to reach a formal legal determination, finding specifically that Brady had collected substantial sums, reported in various accounts as several million euros, from licensees and preorder customers for a product that was never delivered in working form, a resolution that addresses the financial conduct directly rather than adjudicating the underlying physics claim, though the absence of any delivered working unit across the relevant period is itself consistent with the physics-based objection above rather than requiring a separate explanation.

\section{The Entropic Gravity Framing}

The Perendev motor makes no gravitational claim, concerning energy conservation via permanent magnets alone, and is included in this series as a comparative case in the over-unity genre alongside the Yildiz Motor and Joe Newman's device, illustrating how the underlying physics objection, permanent magnets doing no net work over a cycle, recurs with only superficial variation in claimed mechanism across multiple independently promoted devices spanning several decades.

\section{Objections}

Some supporters have argued that the fraud conviction addressed only Brady's business conduct and left open whether a corrected or improved version of the underlying magnetic timing concept might work, distinguishing the legal outcome from a definitive scientific refutation. This is true in the narrow sense that a fraud conviction is not itself a physics paper, but the underlying energy conservation objection applies to any static permanent magnet arrangement regardless of the specific timing or shielding scheme proposed, and no subsequent independently verified working version of the Perendev concept, corrected or otherwise, has been demonstrated in the years since.

\section{Conclusion}

The Perendev motor fails the same permanent-magnet energy conservation test as other over-unity magnet motor claims, and its financing structure, substantial advance revenue collected ahead of independent verification, supplies an additional, physics-independent reason the claim continued to be promoted well past the point any working product should have been demonstrable, culminating in a fraud conviction that addressed the financial conduct without needing to separately adjudicate a physics claim the device's own absence of a delivered product had already left unsupported.

\end{document}
